\section{Scope and Two-Level Scheduling Architecture}
\label{sec:order_scheduling_scope}

The warehouse workload is organized into orders, each consisting of one or
more missions. Every mission carries an order identifier, and all missions
of the same order must be executed by the same operator. In the initial
model, operator, mobile resource, and trolley denote the same scheduling
entity. Mission priority is meaningful only within its order: a larger
\texttt{NM\_PRIORITY} value indicates greater urgency.

The design separates two problems. Offline planning minimizes the number of
resources required to execute the known workload within a shift, respecting
compatibility, order ownership, mission priorities, and order execution-time
limits. Online rescheduling keeps the selected resources fixed and changes
the sequencing of orders to compensate for execution deviations. Its aim is
to reduce deviations from the offline performance reference, including
resource utilization and the fraction of completed orders.

This chapter specifies the proposed architecture, not an already validated
implementation. The earlier mission-wise resource selector and accumulated
workload predictor are a starting point, but do not solve the new order-level
sequencing problem. The following design replaces that formulation as the
target for subsequent implementation.

\section{Order Constraints and Priorities}

Let $\mathcal M$ denote the shift missions, $\mathcal O$ the orders, and
$\mathcal M_o$ the missions of order $o$. Order identifiers define a partition
of $\mathcal M$. The actual input column and its uniqueness scope still
require verification: mission identifiers must not be used as order
identifiers without checking the data.

\subsection{One Operator per Order}

Let $\mathcal R$ be the available resources, $a_{or}$ binary order assignments,
and $z_r$ binary resource-activation variables. Offline planning requires
\begin{align}
 \sum_{r\in\mathcal R}a_{or}&=1 &&\forall o\in\mathcal O,\\
 a_{or}&\leq z_r &&\forall o,r,\\
 a_{or}&=0 &&\forall r\notin\mathcal F_o,
\end{align}
where $\mathcal F_o=\bigcap_{i\in\mathcal M_o}\mathcal F_i$ and
$\mathcal F_i$ contains the operators compatible with mission $i$.
An empty intersection makes the order incompatible with the available
resources under the single-operator requirement.

Every mission of order $o$ inherits its operator $r_o$. This identity persists
until the last mission completes, including across batch boundaries and idle
periods. The proposed online baseline retains the offline order assignments,
so sequencing is its only decision variable. Reassigning an unstarted order
would introduce an additional degree of freedom and is outside this baseline.

\subsection{Priority Within an Order}

Let $p_i$ be the \texttt{NM\_PRIORITY} value of mission $i$. If
$i,j\in\mathcal M_o$ and $p_i>p_j$, mission $i$ is more urgent than $j$.
Values from different orders do not define a global priority ranking. No
order-level priority is inferred from their maximum or average.

Two interpretations must be distinguished. A dispatch preference selects
the higher-priority mission among executable missions of the same order.
A mandatory precedence instead requires
\begin{equation}
 f_i\leq s_j,
\end{equation}
where $s_j$ and $f_i$ are start and completion times: mission $j$ must then
wait even when mission $i$ is temporarily unavailable. Neither interpretation
creates precedence between unrelated orders. Priorities 9 and 4 in order A,
for example, do not prevent a mission of order B with priority 2 from
starting on another operator.

The priority direction and within-order scope are established. Dispatch
preference versus mandatory precedence, and tie-breaking for equal values,
remain to be specified. Offline and online scheduling must use the same
explicit rule. Technological precedences remain hard constraints regardless
of the chosen priority interpretation.

\subsection{Deadline as Maximum Elapsed Execution Time}

An optional parameter $D_o>0$ is entered manually for order $o$; it is not
currently supplied by the input files. It denotes a maximum elapsed execution
duration, rather than a due date measured from shift start. Define
\begin{equation}
 S_o=\min_{i\in\mathcal M_o}s_i,\qquad
 C_o=\max_{i\in\mathcal M_o}f_i.
\end{equation}
For every order with a specified limit, the requirement is
\begin{equation}
 C_o-S_o\leq D_o.
 \label{eq:order_duration_limit}
\end{equation}
The clock starts at the actual start of the first mission and ends at the
actual completion of the last. Waiting between missions, including time
spent serving another order, counts toward this duration. Waiting before
the first mission does not. Therefore this constraint alone does not penalize
postponing an order's start; shift completion and progress against the offline
reference must also be considered.

Once an order starts, its absolute expiry time is
\begin{equation}
 d_o=S_o^{\mathrm{actual}}+D_o.
\end{equation}
The observed start and expiry are retained at subsequent updates. For an
unstarted order, each candidate checks its predicted duration
$\widehat C_o-\widehat S_o\leq D_o$. Entering a new batch never resets the
clock. Orders without $D_o$ have no additional order-duration constraint,
but remain part of the shift workload.

The project PDF discusses estimating a duration from mission count and mean
mission time. This can support the manual choice of $D_o$, but does not
replace the required definition or provide a calibrated value. Nominal
duration components, including travel and handling, must be specified before
using such estimates.

\section{Offline Resource Dimensioning and Reference Plan}

For a shift duration $T_{\mathrm{shift}}$, the conceptual offline problem is
\begin{equation}
 \min_{z,a,\sigma,s,f}\sum_{r\in\mathcal R}z_r,
\end{equation}
subject to the assignment conditions, compatibility, the agreed priority
rule, feasible routes, resource non-overlap, and
\begin{equation}
 0\leq s_i<f_i\leq T_{\mathrm{shift}}\quad\forall i\in\mathcal M,
 \qquad C_o-S_o\leq D_o
\end{equation}
for orders with a duration limit. Here $\sigma$ describes execution sequences.
Durations can depend on the resource and its predecessor through repositioning
travel. Among plans with the same resource count, completion time and load
balance can be secondary criteria.

A heuristic returns the smallest feasible resource configuration found,
unless minimality is certified. Failure to find a schedule is not proof of
infeasibility. The shift duration is an input constraint: a soft target around
eight hours does not automatically authorize overtime.

The selected plan fixes $\mathcal R^*$, $q^*=|\mathcal R^*|$, and $r_o$.
It provides reference completion curves at physical time $t$:
\begin{align}
 N^{\mathrm{ref}}(t)&=\sum_{i\in\mathcal M}\mathbf1\{\bar f_i\leq t\},\\
 O^{\mathrm{ref}}(t)&=\sum_{o\in\mathcal O}\mathbf1\{\bar C_o\leq t\}.
\end{align}
Resource utilization references come from active working times in this same
plan. The reference obeys the online operational constraints and is preserved
across updates rather than reset to conceal execution deviations.

\section{Batch Clock and Receding Horizon}

The baseline uses batches of fixed mission cardinality $m$, except for a
possibly smaller final batch. At step $k$, batch $k$ has just been executed.
The controller predicts the next $H$ batches and applies only the schedule
for batch $k+1$. With $H=3$, it evaluates batches $k+1$, $k+2$, and $k+3$,
executes the first, observes its outcome, and repeats from the updated state.

For a concrete baseline, the offline mission list is partitioned into ordered
batches whose membership is retained online. The controller changes the order
of service within upcoming batches subject to ownership and intra-order
rules. This limits the decision set to at most $Hm$ upcoming missions instead
of the full residual workload. Moving missions between batches would require
an additional admission rule and is not assumed here. Fixed membership is a
baseline design choice whose restriction on achievable improvements must be
evaluated.

An order may cross several batches, retaining its operator, first-start time,
priorities, and duration limit. Crossing a boundary neither closes nor
restarts the order. Whether different orders may be interleaved on one
operator is a separate policy still to be settled.

Let $t_k$ be the physical completion time of batch $k$. The baseline uses a
global batch-completion event, so the elapsed interval
\begin{equation}
 h_{k+1}=t_{k+1}-t_k
\end{equation}
is variable and includes the time until the last resource finishes its batch
work. Earlier-finishing resources may wait at the boundary; this overhead
must be measured. Each candidate has its own predicted batch-completion times
and is compared with the reference at those same physical times.

The project PDF also proposes 30-minute periodic windows and adaptive
mini-batch sizes. These are alternative designs, not equivalent to fixed-size
batches: three batches do not necessarily span 90 minutes. They remain
possible extensions. The baseline follows the batch-indexed formulation;
$H=3$ is an initial configuration to evaluate, not a validated computational
choice. The number $n$ of candidate schedules differs from the horizon $H$.

\section{Persistent State and Predictive Model}

Accumulated workload alone does not represent order ownership, elapsed
deadlines, and execution progress. A proposed state is
\begin{equation}
\begin{split}
 x_k=\big(&t_k,\mu_k,\mathcal W_k,\mathcal I_k,\mathcal C_k,
 \{p_r,b_r,\rho_r,A_r\}_{r\in\mathcal R^*},\\
 &\{r_o,S_o,d_o\}_{o\in\mathcal O},\widehat\theta_k\big).
\end{split}
\end{equation}
Here $\mu_k$ is the CPN marking; $\mathcal W_k$, $\mathcal I_k$, and
$\mathcal C_k$ are unstarted, active, and completed missions. Resource fields
are position $p_r$, availability $b_r$, estimated residual processing time
$\rho_r$, and accumulated active working time $A_r$. Start and expiry fields
are unset until applicable; $\widehat\theta_k$ contains duration-model
estimates. Mission records retain order identifiers, priorities, and batch
membership. Pending events and clocks must also be retained for resumption
inside an execution.

For a fixed workload, conservation requires
\begin{equation}
 \mathcal M=\mathcal W_k\mathbin{\dot\cup}\mathcal I_k
 \mathbin{\dot\cup}\mathcal C_k.
\end{equation}
At normal batch-completion updates, the active set is empty under the global
barrier convention. It remains necessary for simulation events and possible
future periodic updates.

The input $u_{k+1}$ specifies order sequences for the next batch on the fixed
resources. Intra-order choices follow the agreed priority rule. The dynamics
are represented by an event-driven transition:
\begin{align}
 x_{k+1}&=F(x_k,u_{k+1},w_{k+1}),\\
 \widehat x_{k|k}&=x_k,\\
 \widehat x_{k+\ell|k}&=\widehat F(\widehat x_{k+\ell-1|k},
 u_{k+\ell|k};\widehat\theta_k),\quad\ell=1,\ldots,H.
\end{align}
The actual disturbance $w_{k+1}$ is unknown at planning time. The predictor
uses available observations and estimates, never future realized durations.
A workload update $T_r^+=T_r+\Delta T_r$ remains useful as accounting, but a
constant-matrix workload model is not a complete sequencing model.

\section{CPN and Path Feasibility}

The CPN enforces admissible starts, compatibility, and resource availability.
An order register or equivalent persistent token data enforces $r_o$ across
missions. Reserving the resource only during a single mission is insufficient.
For execution intervals $[s_i,f_i)$, capacity requires
\begin{equation}
 \sum_{i:r_{o(i)}=r}\mathbf1\{s_i\leq t<f_i\}\leq1\qquad\forall r,t.
\end{equation}
Starts are controllable scheduling decisions; completions follow from the
execution dynamics and cannot be freely selected by the optimizer.

A* supplies feasible paths and travel estimates from predicted resource
positions. Sequence changes can alter repositioning times and must therefore
be reflected in prediction. Candidates violating compatibility, capacity,
technological precedences, or the agreed priority rule are rejected. If the
assigned operator cannot reach a mission, that mission cannot simply be
given to another operator. Deferral preserves its identity and order history
and cannot be reported as batch completion; an explicit recovery is needed.

\section{Performance Outputs and Online Objective}

Observed counts $N(t)$ and $O(t)$ use actual completion events. An order is
complete only after all its missions finish. For a nonempty order set,
\begin{equation}
 c_O(t)=\frac{O(t)}{|\mathcal O|},\qquad
 c_O^{\mathrm{ref}}(t)=\frac{O^{\mathrm{ref}}(t)}{|\mathcal O|}.
\end{equation}
Multiplying by 100 gives the completed-order percentage. Resource utilization
for $t>0$ is
\begin{equation}
 \eta_r(t)=\frac{A_r(t)}{t}.
\end{equation}
The denominator is elapsed time since shift start; scheduled breaks would
require a consistent adjustment in reference and observation. Active work
excludes waiting. Utilization at zero elapsed time is not used in the cost.

The signed mission error is $e_N(t)=N^{\mathrm{ref}}(t)-N(t)$. A negative
value denotes progress ahead of the reference. It is already included in the
cumulative counts and must not be added again as a separate credit. Extra
short missions do not establish enough available time to complete a long
order within its duration limit.

For $U_k=(u_{k+1|k},\ldots,u_{k+H|k})$, a proposed dimensionless cost is
\begin{align}
 J_H={}&\sum_{\ell=1}^H\Bigg(
 w_O\left[c_O^{\mathrm{ref}}(\widehat t_{k+\ell|k})
 -\widehat c_{O,k+\ell|k}\right]_+^2\nonumber\\
 &+\frac{w_\eta}{q^*}\sum_{r\in\mathcal R^*}
 \left[\widehat\eta_{r,k+\ell|k}
 -\eta_r^{\mathrm{ref}}(\widehat t_{k+\ell|k})\right]^2\Bigg)
 +V_f(\widehat x_{k+H|k}),
\end{align}
where $[v]_+=\max(0,v)$ and weights are nonnegative and require tuning.
The completion term penalizes falling behind without penalizing early
completion. Utilization tracking is subordinate to feasibility and must be
checked for unwanted incentives to create idle time.

One possible terminal term is
\begin{equation}
 V_f(x)=w_W\frac{\widehat W(x)}{W_s},\qquad W_s>0,
\end{equation}
where $\widehat W$ estimates remaining work, including active residual times.
It discourages ignoring residual work but is not an exact completion-time
estimate or a stability guarantee. With fixed membership and sequence-
independent durations, it may be equal across candidates and then cannot
distinguish their sequences.

For orders continuing beyond the prediction horizon, the controller must
retain elapsed clocks and evaluate a feasible continuation to completion,
or explicitly flag that full deadline feasibility has not been established.
A lower bound on remaining duration detects some impossible cases, but
passing it does not certify feasibility. These checks can inspect residual
order metadata beyond the $Hm$ decision missions without optimizing the
entire remaining workload. An unfinished order is never assigned zero
lateness risk simply because its completion lies outside the horizon.

The nominal baseline treats specified maximum durations as constraints.
Replacing them with lateness penalties is a different service policy still
to be agreed for recovery. If adopted explicitly, lateness is
$[C_o-S_o-D_o]_+$; its weight is not inferred by comparing mission priorities
across orders.

\section{Online Optimization and Execution Loop}

The finite-horizon scheduling problem is
\begin{equation}
 U_k^*\in\underset{U_k\in\mathcal U_H(x_k)}{\operatorname{arg\,min}}
 J_H(x_k,U_k),\qquad u_{k+1}=u_{k+1|k}^*.
\end{equation}
The feasible set includes fixed resources and order assignments, batch
membership, priority rules, paths, capacity, shift limits, and applicable
order-duration constraints. Near the end of the list, the predicted batch
count is reduced to the remaining number.

A first implementation can use bounded simulation-based search. Each
candidate starts from an independent copy of the observed state and is
evaluated over $H$ batches. Order swaps and insertions are possible sequence
moves, admitted only when consistent with committed constraints. A common
continuation heuristic can limit search in future batches. The result is
then the best feasible candidate found, not a certified global optimum.

The loop is:
\begin{enumerate}
 \item Observe the completed batch and update completion events, resource
 states, order clocks, and model estimates.
 \item Retain the next $H$ batches and reconstruct a baseline candidate from
 the previous plan, rechecking it against the observed state.
 \item Generate alternative order sequences within a computation budget
 $\tau_{\mathrm{calc}}$ and simulate each from the same initial state.
 \item Verify constraints, including obligations of orders crossing the
 horizon, and evaluate the performance cost.
 \item Keep the best verified candidate and apply only its first batch;
 the remaining predicted batches are provisional.
 \item Execute with actual disturbances, record events, and repeat after
 batch completion. Stop when the workload is complete or the shift-end
 and recovery policy requires stopping.
\end{enumerate}

Computation time includes candidate generation, path evaluation, simulation,
and state reconciliation. Latency must be included in the execution timeline
or explicitly assumed negligible in an initial experiment. Batch size $m$,
horizon $H$, candidate count, and budget require measured calibration.
Smaller batches produce more frequent feedback but more updates and possible
synchronization overhead.

When no verified candidate is found, preserve ownership and existing
commitments and report the cause. Resource unavailability, an expired
deadline, budget exhaustion, and proven infeasibility are distinct outcomes.
Constraints must not be silently relaxed to return a schedule.

\section{Implementation Decisions and Validation}

The outstanding decisions from the project document are updated as follows:
\begin{enumerate}
 \item \textbf{Order identifier:} verify the column and uniqueness scope.
 \item \textbf{Priority direction:} larger \texttt{NM\_PRIORITY} means
 greater urgency within its order; tie-breaking remains open.
 \item \textbf{Priority enforcement:} choose dispatch preference or mandatory
 intra-order precedence using the distinction explained above.
 \item \textbf{Order continuity:} crossing batches is allowed; interleaving
 different orders on one operator remains to be decided.
 \item \textbf{Resource identity:} operator and trolley coincide initially.
 \item \textbf{Deadline origin:} actual first-mission start, as required by
 the elapsed-duration definition. Manual values and nominal duration
 components still need specification.
 \item \textbf{Deadline recovery:} limits are nominal constraints; decide
 how to handle unavoidable violations during execution.
 \item \textbf{Prediction and search:} test $H=3$ initially; calibrate $m$,
 budget, and continuation. Fixed batch membership and offline ownership
 are explicit baseline choices to validate.
 \item \textbf{Update timing:} use batch completion, without a mandatory
 30-minute wait. Periodic or asynchronous updates require another policy.
\end{enumerate}

Validation starts with small deterministic examples: single-mission orders,
orders crossing batches, different priorities within one order, parallel
orders, missing optional deadlines, incompatible orders, and a final partial
batch. Check for lost or duplicated missions, operator changes, invalid
overlaps, reset clocks, and premature completion records. Interrupted and
resumed simulation must match continuous execution under fixed decisions.

Compare the unchanged offline plan, a one-batch controller, and the
multi-batch controller with identical resources, assignments, workload, and
priority policy. Disturbances must be comparable across policies and hidden
from the predictor until observed. Report completed-order fraction, mission
progress, utilization, duration violations, work remaining at shift end,
batch waiting, and computation times. Incomplete orders remain in service
metric denominators.

Finding a feasible candidate does not prove recursive feasibility under
disturbances or eventual completion. These require further assumptions and
analysis. No experimental improvement is claimed before implementation and
validation of the proposed architecture.

